\chapter{Arrays, tables, and visual evidence}
\label{ch:arrays-tables-viz}

\begin{learningobjectives}
After this chapter, you should be able to:
\begin{itemize}
  \item reason about a NumPy array using shape, axes, and data type;
  \item distinguish vectorization from a mere syntax shortcut;
  \item use pandas labels while guarding alignment and missingness; and
  \item design a visualization as an honest, reproducible scientific argument.
\end{itemize}
\end{learningobjectives}

\begin{chapterconnection}
The preceding chapters built trustworthy Python interfaces. We now use mature
scientific packages through those interfaces. NumPy adds regular numerical
structure, pandas adds labels and heterogeneous columns, and visualization
turns selected relationships into arguments that other people must interpret.
\end{chapterconnection}

\section{From Python containers to arrays}

A Python list can hold references to heterogeneous objects. A NumPy
\term{array} stores a regular, multidimensional collection governed by a shape
and data type. That structure supports compact memory layouts and operations
implemented over entire blocks of values.

For an array with shape $(m,n)$, axis 0 indexes $m$ rows and axis 1 indexes $n$
columns. An operation such as \code{mean(axis=0)} reduces rows and returns one
mean per column. ``Across rows'' is ambiguous; name the axis and resulting
shape.

\begin{lstlisting}
import numpy as np

measurements = np.array(
    [[20.1, 101.2], [20.4, 100.9], [20.8, 101.0]],
    dtype=np.float64,
)

column_means = measurements.mean(axis=0)
assert measurements.shape == (3, 2)
assert column_means.shape == (2,)
\end{lstlisting}

The data type is part of the numerical contract. Integer overflow, floating
roundoff, missing-value representation, and implicit coercion can change a
result without changing array shape.

\section{Vectorization and broadcasting}

Vectorization expresses operations over arrays so optimized compiled loops can
perform the elementwise work. It often improves clarity and performance, but it
does not remove algorithmic cost or memory allocation.

Broadcasting compares dimensions from the right and permits compatible
singleton dimensions to act as if repeated. Always predict the result shape
before relying on broadcasting. An operation that runs can still combine the
wrong scientific axes.

\begin{warningbox}
Vectorization can allocate enormous temporary arrays. Measure the working set,
use in-place operations only when ownership is safe, and consider chunked or
streaming computation when the full result cannot fit comfortably in memory.
\end{warningbox}

\section{Shape is syntax; axis meaning is semantics}

NumPy can verify whether dimensions are compatible. It cannot know what the
dimensions mean. An array with shape $(100,3)$ might represent 100 patients and
three biomarkers, 100 time points and three spatial coordinates, or 100
simulations and three model parameters. The same matrix multiplication can be
dimensionally valid in all three cases while answering entirely different
questions.

Before an important operation, write a small shape ledger:

\begin{center}
\begin{tabular}{@{}lll@{}}
\toprule
\textbf{Name} & \textbf{Shape} & \textbf{Axis meaning} \\
\midrule
\code{x} & $(n,d)$ & observations $\times$ features \\
\code{beta} & $(d,)$ & one coefficient per feature \\
\code{x @ beta} & $(n,)$ & one score per observation \\
\bottomrule
\end{tabular}
\end{center}

The ledger makes the intended contraction visible: the feature axis of
\code{x} meets the feature axis of \code{beta}. If \code{x} were transposed,
the operation should fail rather than be reshaped until it runs. A reshape
changes an interpretation unless you can explain where every axis went.

\begin{definitionbox}{vectorization}
Vectorization expresses a collection of related operations through an array
interface whose looping occurs in optimized lower-level code. It changes where
iteration is expressed; it does not make the iteration, memory traffic, or
algorithmic complexity disappear.
\end{definitionbox}

\section{Data type is a numerical design decision}

An array's \term{dtype} determines how each value is represented and which
operations are available. Integers are exact inside a finite range. Floating
point values cover a much larger dynamic range but approximate most real
numbers. Boolean arrays support logical masks. Date-time types attach a time
unit. Strings and generic Python-object arrays have different storage and
performance behavior from regular numeric arrays.

Three questions belong in every numerical review:

\begin{enumerate}
  \item \textbf{Range:} can the representation hold the largest magnitude that
  an intermediate computation may produce?
  \item \textbf{Precision:} can it distinguish changes that matter to the
  analysis?
  \item \textbf{Missingness:} how will absent, invalid, censored, or not-yet-
  observed values be represented without confusing their meanings?
\end{enumerate}

For floating point work, equality is often the wrong contract. A computed value
may differ from an algebraically equivalent result by a small rounding error.
Use a tolerance derived from the problem's scale and numerical method, then
state whether the tolerance is absolute, relative, or both. The convenience of
\code{np.allclose} does not choose a scientifically meaningful tolerance.

For sums of values with very different magnitudes, operation order may matter.
For probabilities near zero or one, algebraically equivalent formulas can have
very different numerical stability. Numerical analysis is not a repair applied
after modeling; it is part of deciding what computation the model actually
performs.

\section{pandas adds labeled structure}

A pandas \term{Series} combines values with an index. A \term{DataFrame}
combines columns, row labels, and heterogeneous column data types. Labels reduce
some positional errors and introduce alignment behavior that must be understood.

Arithmetic between Series aligns by index label, not necessarily by physical
position. A join encodes assumptions about key uniqueness and relationship
cardinality. Missing values interact with each column's data type and with the
semantics of an aggregation.

\begin{professionalbox}
Before a groupby, join, or pivot, state what one row represents, which keys are
unique, whether missing groups should appear, and how row counts should change.
Check those claims after the operation.
\end{professionalbox}

Tables are not automatically tidy, independent, or statistically valid. pandas
makes transformations expressive; scientific meaning still belongs to the
analyst and data contract.

\subsection*{State what one row means}

Suppose a table has columns \code{experiment}, \code{sensor}, \code{time\_s},
and \code{temperature\_c}. If one row is one sensor reading, then grouping by
experiment and taking a mean produces one row per experiment only after we
decide how sensors and time points should be weighted. If one sensor recorded
twice as often, a naive mean weights that sensor twice as much. pandas will
execute the aggregation faithfully; it cannot determine whether that weighting
matches the estimand.

Indexes also carry assumptions. A default integer index may merely identify
current row positions. A domain index such as \code{experiment\_id} claims that
the labels have stable meaning and perhaps uniqueness. Set an index because its
semantics support an operation, not because it makes the display attractive.

\begin{definitionbox}{tidy observational table}
In a common tidy representation, each row is one observational unit, each
column is one variable, and each cell contains one value. This convention is a
useful design tool, not a proof that observations are independent or that the
table is suitable for a particular statistical model.
\end{definitionbox}

\section{Visualization is a transformation}

A figure maps data and decisions into marks, positions, color, size, facets,
scales, labels, and annotations. Every choice emphasizes some comparisons and
suppresses others.

Matplotlib supplies low-level control and a broad scientific ecosystem. Seaborn
adds statistical plotting conventions over pandas data. Plotly supports
interactive graphics. Tool choice follows the communication goal, audience,
output medium, accessibility requirements, and reproducibility needs.

An honest workflow records:

\begin{itemize}
  \item the data snapshot and filtering rules;
  \item units, transformations, aggregation, and uncertainty;
  \item scale limits and whether an axis is linear or logarithmic;
  \item random seed and simulation parameters;
  \item color choices that remain interpretable without color alone; and
  \item the code and package versions that produced the artifact.
\end{itemize}

\subsection*{Choose an encoding from the comparison}

A chart type should follow the question rather than habit:

\begin{longtable}{@{}>{\raggedright\arraybackslash}p{0.28\textwidth}
>{\raggedright\arraybackslash}p{0.28\textwidth}
>{\raggedright\arraybackslash}p{0.35\textwidth}@{}}
\toprule
\textbf{Question} & \textbf{Useful starting point} & \textbf{Common risk} \\
\midrule
\endhead
How does one quantity vary with another? & scatter plot, perhaps with a
model summary & overplotting, hidden groups, implying causation \\
How does a quantity evolve in ordered time? & line plot with observed points &
connecting missing intervals or unrelated units \\
How do distributions differ across groups? & ECDF, interval plot, box or violin
plot with observations & hiding sample size, multimodality, or unequal support \\
How uncertain is an estimate? & point and interval with the estimand named &
unlabeled interval type or visually treating uncertainty as error \\
How are values arranged over two dimensions? & heat map with a justified scale
and color map & perceptual distortion, inaccessible color, hidden normalization \\
\bottomrule
\end{longtable}

Position along a common scale usually supports more precise comparison than
area, volume, or angle. Color can reveal grouping or magnitude, but it should
not carry the only copy of essential information. Interactive hover text can
add detail; the unhovered figure still needs a title, axes, units, legend,
caption, and visible central claim.

\section{Simulation as executable reasoning}

Simulation can make probability and dynamical behavior visible. A random seed
supports reproducibility of one run; it does not prove robustness. Compare
multiple seeds or summarize a distribution when variability matters.

For a Monte Carlo estimator $\widehat{\theta}_n$, distinguish numerical error,
sampling variability, model assumptions, and implementation defects. More
samples may reduce Monte Carlo variance while leaving structural bias intact.

\begin{workedexample}{use probability theory to test a simulation}
Let $X_1,\ldots,X_n$ be independent Bernoulli random variables with success
probability $p$. The sample mean

\[
  \widehat{p}=\frac{1}{n}\sum_{i=1}^{n}X_i
\]

has expectation $p$ and standard deviation
$\sqrt{p(1-p)/n}$. We can use this mathematical result as an independent oracle
for a vectorized simulation.

\begin{lstlisting}
import numpy as np

rng = np.random.default_rng(438)
p = 0.30
n = 1_000
repetitions = 5_000

success_counts = rng.binomial(n=n, p=p, size=repetitions)
estimates = success_counts / n

empirical_center = estimates.mean()
empirical_sd = estimates.std(ddof=1)
theoretical_sd = np.sqrt(p * (1 - p) / n)

assert estimates.shape == (repetitions,)
assert np.isclose(empirical_center, p, atol=0.01)
assert np.isclose(empirical_sd, theoretical_sd, rtol=0.05)
\end{lstlisting}

\textbf{Step 1: identify the simulated object.} Each entry of
\code{success\_counts} is one binomial count summarizing 1,000 Bernoulli trials.
The array axis indexes repeated experiments, not individual observations.

\textbf{Step 2: vectorize at the right level.} We request 5,000 independent
counts directly. Materializing a $(5000,1000)$ Boolean array would also work,
but it would consume more memory to answer the same probability question.

\textbf{Step 3: separate reproducibility from validity.} The local random-number
generator and seed reproduce this particular stream. They do not establish that
the binomial model is appropriate for real observations.

\textbf{Step 4: compare with theory.} The analytic standard deviation is a test
oracle independent of the random sample. Tolerances reflect expected Monte
Carlo fluctuation; they are not selected merely to make the assertions pass.

\textbf{Step 5: visualize the sampling distribution.} A histogram or empirical
cumulative distribution function should be labeled as a distribution of
\emph{estimates across repeated experiments}. It is not the distribution of
individual Bernoulli observations.
\end{workedexample}

This example illustrates a reusable scientific pattern: derive what should be
true, implement the simulation, compare independent evidence, and explain the
remaining modeling assumptions. Simulation complements probability theory; it
does not replace it.

\begin{checkpointbox}
You receive a table with one row per sensor reading. Design a figure comparing
two instruments over time. State the row meaning, units, missing-data policy,
aggregation, uncertainty representation, accessibility choices, and two tests
for the transformation feeding the plot.
\end{checkpointbox}

\section{Worked example: standardize three sensor channels}

Suppose rows represent time points and columns represent sensors. We want to
subtract each sensor's mean and divide by its standard deviation. The formula
for entry $(i,j)$ is

\[
  z_{ij}=\frac{x_{ij}-\bar{x}_j}{s_j}.
\]

The subscript $j$ matters: statistics are computed down rows separately for
each column.

\begin{workedexample}{predict shape before execution}
\begin{lstlisting}
import numpy as np

x = np.array(
    [
        [10.0, 100.0, 1.1],
        [12.0, 104.0, 0.9],
        [14.0, 102.0, 1.0],
        [16.0, 106.0, 1.2],
    ]
)

means = x.mean(axis=0)
scales = x.std(axis=0, ddof=1)
z = (x - means) / scales
\end{lstlisting}

\textbf{Step 1: annotate shapes.} \code{x.shape} is $(4,3)$. Reducing axis 0
gives \code{means.shape == (3,)} and \code{scales.shape == (3,)}.

\textbf{Step 2: predict broadcasting.} NumPy compares trailing dimensions.
The length-3 vectors align with the three columns and act across four rows.
The result retains shape $(4,3)$.

\textbf{Step 3: check invariants.} Each standardized column should have mean
close to zero and sample standard deviation close to one.

\begin{lstlisting}
assert z.shape == x.shape
assert np.allclose(z.mean(axis=0), 0.0)
assert np.allclose(z.std(axis=0, ddof=1), 1.0)
\end{lstlisting}

\textbf{Step 4: ask scientific questions.} Is standardization appropriate for
all three sensors? Are the observations independent? Should calibration ranges
or robust statistics replace the sample mean and standard deviation? Shape
correctness does not answer these questions.
\end{workedexample}

If we mistakenly reduce \code{axis=1}, NumPy may still produce values after a
reshape. The program can run while standardizing each time point across sensors
with incompatible units. Predicting dimensions and naming their meaning catches
an error that syntax cannot.

\section{Worked example: pandas alignment is not row order}

Consider two Series of corrections indexed by sensor identifier. pandas aligns
labels before arithmetic.

\begin{lstlisting}
import pandas as pd

reading = pd.Series({"A": 10.0, "B": 20.0})
offset = pd.Series({"B": 1.5, "A": 0.5})
corrected = reading - offset

assert corrected["A"] == 9.5
assert corrected["B"] == 18.5
\end{lstlisting}

The physical order differs, but labels preserve the intended pairing. If one
Series uses \code{"sensor-A"} and the other uses \code{"A"}, alignment produces
missing values rather than guessing equivalence. That behavior is safer than
silent positional pairing only when the analyst checks the resulting index and
missingness.

\begin{misconceptionbox}
\textbf{Misconception:} pandas prevents incorrect joins because columns have
names.

\textbf{Correction:} names permit explicit alignment; they do not prove key
uniqueness, cardinality, unit compatibility, or semantic identity. A many-to-many
join can multiply rows exactly as requested and still invalidate an analysis.
\end{misconceptionbox}

\section{From analytical question to figure}

Suppose the question is whether sensor bias changes with temperature. A useful
figure might place reference temperature on the horizontal axis and residual
on the vertical axis, show individual observations with transparency, add a
zero reference line, and facet by sensor model. A line chart connecting
unrelated samples would imply an ordering that does not exist.

Before plotting, write the sentence the figure should support. Then identify
the comparison, encodings, uncertainty, and possible confounders. After
plotting, try to state a false conclusion the design might encourage. This
adversarial reading is part of visual quality assurance.

\conceptcheck{Why might truncating the residual axis make a small bias appear
important? When is a nonzero baseline nevertheless appropriate?}

\section{Exercises}

\subsection*{Foundational}
\begin{enumerate}
  \item For shapes $(20,4)$ and $(4,)$, explain broadcasting. Repeat for
  $(20,4)$ and $(20,)$. Which operation fails, and why?
  \item Distinguish NumPy data type, array shape, and axis meaning.
  \item Explain why a random seed makes one simulation reproducible but does not
  establish that its conclusion is stable.
  \item For a Bernoulli sample mean, derive how the theoretical standard
  deviation changes when $n$ is multiplied by four.
\end{enumerate}

\subsection*{Applied}
\begin{enumerate}
  \item Modify the standardization example to retain means with shape $(1,3)$
  using \code{keepdims=True}. Explain whether this changes values or only makes
  broadcasting intent more visible.
  \item Construct two pandas Series whose inconsistent labels produce missing
  output. Add assertions that detect the mismatch before arithmetic.
  \item Modify the binomial experiment across four values of $n$. Compare the
  observed standard deviations with theory and design one figure that makes the
  $1/\sqrt{n}$ relationship visible.
  \item Design a static and interactive version of the same scientific figure.
  State what interaction adds and what information must remain visible without
  hovering.
\end{enumerate}

\subsection*{Design}
Write a figure specification before writing plotting code. Include audience,
claim, data snapshot, row meaning, units, encodings, aggregation, uncertainty,
accessibility, caption, and two misleading alternatives you rejected.

\begin{chapterreview}
\textbf{Key ideas.} NumPy arrays require joint reasoning about values, shape,
axes, and dtype. Broadcasting is a dimensional rule, not evidence of semantic
compatibility. pandas aligns labels and therefore requires explicit key and
missingness checks. A visualization is a reproducible transformation and
argument, not decoration.

\textbf{Vocabulary.} array; shape; axis; dtype; vectorization; broadcasting;
Series; DataFrame; index; alignment; join cardinality; scale; encoding;
uncertainty; simulation; Monte Carlo error.

\textbf{Explain aloud.} How can an array operation have the correct output shape
and still combine the wrong scientific quantities?
\end{chapterreview}

\begin{notebooklink}
Lecture 3 notebooks 03 and 04 connect native Python containers to NumPy and
pandas, then develop static, statistical, interactive, and simulation-based
visualizations. Reproduce the probability experiment there before changing its
sample size, seed, or visualization.
\end{notebooklink}

\section{Takeaway}

Arrays make structure computationally explicit, tables add labels, and figures
turn selected relationships into visible evidence. None of these layers chooses
the scientific contract automatically.
